\section{Six Birds: how the primitives specialize to geometry}
\label{sec:six-birds}

Six Birds Theory~\cite{SixBirdsTheory} describes how a higher layer of description (a ``theory'' in its technical sense: a lens, a completion and an audit) becomes available when repeated compression of a lower layer stabilizes.
It identifies six primitives.
The primitives are unchanged here; what changes is the object they produce.
For counting, the stable object is quantity~\cite{ToCountAStone}; for plotting, it is location, distance and transport.
This section states the correspondence informally; later sections make each item precise.

\paragraph{P1, operator rewrite (closure).}
A layer is usable only if its dynamics can be written in its own variables without reopening micro detail.
Here this is the macro kernel $K=UP^{\tau}C$ and the closure operator $E=P^{\tau}CU$ on micro distributions.
Closure is not assumed; its defect is measured (\cref{sec:closure}).

\paragraph{P2, constraints (feasibility).}
Constraints decide which transitions are possible and how likely they are.
Geometrically they shape neighbourhoods and anisotropy.
Gating moves in one direction deforms the induced metric even when the lens is held fixed (\cref{sec:anisotropy}).

\paragraph{P3, protocols (order of composition).}
A protocol is a chain of moves.
Distance is the cheapest protocol between two points, and curvature is the dependence of transport on the route taken: carrying a frame around a loop need not return it unchanged.
\Cref{sec:curvature} shows that this route dependence is exactly the failure of two covariant transport operators to commute, a statement that remains meaningful even though planar rotations themselves commute.

\paragraph{P4, staging (multi-scale refinement).}
Staging supplies the time parameter $\tau$ and a ladder of refining lenses.
Geometry is claimed only for structure that persists along a declared refinement family; in our constructions both the substrate and the lens are refined jointly.

\paragraph{P5, packaging (quotienting indistinguishability).}
Packaging maps many microstates to one macro state.
This is where points come from: a point is a block of microstates that the interface does not distinguish at the chosen resolution.

\paragraph{P6, accounting (cost).}
Accounting assigns a cost to each move.
Our cost is the negative log-probability of a macro transition, so distance is a ledger rather than a ruler.
Different ways of composing the same ledger give different geometries: chained macro transitions give an $L^1$ metric on the grid, while the one-shot cost of a long staged displacement is quadratic and gives Euclidean distance (\cref{sec:pythagoras}).

\begin{table}[t]
\centering
\small
\caption{The six primitives and their role in this paper.}
\label{tab:six-birds-geometry}
\begin{tabularx}{\textwidth}{@{}>{\raggedright\arraybackslash}p{0.2\textwidth}X@{}}
\toprule
Primitive & Role in emergent geometry \\
\midrule
P1 Operator rewrite & Macro kernel and closure operator; closure and persistence defects (\cref{sec:closure}). \\
P2 Constraints & Feasible moves; anisotropy and deformation of the induced metric (\cref{sec:anisotropy}). \\
P3 Protocols & Distance as cheapest protocol; curvature as route dependence of transport (\cref{sec:construction,sec:curvature}). \\
P4 Staging & Stage $\tau$ and refinement families; coherence as a joint limit (\cref{sec:closure,sec:constructed}). \\
P5 Packaging & Points as blocks of microstates; supplied versus learned lenses (\cref{sec:learned,sec:constructed}). \\
P6 Accounting & Negative-log costs; quadratic accounting and the Pythagorean law (\cref{sec:construction,sec:pythagoras}). \\
\bottomrule
\end{tabularx}
\end{table}
